\editionPage[\pdfbookmark[0]{Corrigés des bilans}{bm-bilan-A-cor-auto-parent}]{bilan-A-cor-auto}{Bilan A · Corrigé}{Automatismes expliqués}{\hyperref[bilan-A-auto]{Énoncé p.\,\pageref*{bilan-A-auto}}}
\begin{blocEdition}{bilan-A-cor-auto}
\begin{questions}
\item $9/12-2/12=\res{7/12}$.
\item $\sqrt{25\times3}=\res{5\sqrt3}$.
\item $\res{x^2-8x+16}$. Le double produit vaut $-8x$.
\item La division négative renverse l’ordre : $\res{x\ge-3}$.
\item Centre 1, rayon 2 : $\res{[-1;3]}$.
\item $f(-2)=-6+1=\res{-5}$.
\item Le dénominateur ne doit pas s’annuler : $\res{\mathbb R\setminus\{-5\}}$.
\item Les moyennes de coordonnées donnent $\res{(3;3)}$.
\item $1+0{,}10=\res{1{,}10}$.
\item $1-0{,}35=\res{0{,}65}$.
\end{questions}
\end{blocEdition}
\editionPage{bilan-A-cor-1}{Bilan A · Corrigé 1}{L’aire d’un rectangle}{\hyperref[bilan-A-ex-1]{Énoncé p.\,\pageref*{bilan-A-ex-1}}}
\begin{blocEdition}{bilan-A-cor-1}
\editionEtape{1. Calculer et changer de forme}
$f(0)=0$, $f(2)=8$ et $f(6)=0$. Le développement de $9-(x-3)^2$ donne $6x-x^2=x(6-x)$.
\editionEtape{2. Exploiter un carré}
La fonction croît de 0 à 3 puis décroît de 3 à 6. Son maximum est $\res9$, atteint en $\res{x=3}$.
\editionEtape{3. Résoudre un seuil}
$9-(x-3)^2\ge8$ équivaut à $(x-3)^2\le1$, donc $\res{x\in[2;4]}$.
\editionEtape{4. Interpréter}
Le rectangle optimal est un carré de côté $\res{3\text{ cm}}$, d’aire $\res{9\text{ cm}^2}$ et de périmètre $\res{12\text{ cm}}$. Les valeurs 0 et 6 décrivaient des rectangles aplatis ; le maximum appartient bien au domaine des rectangles non dégénérés.
\editionRetenir{Le contexte ne dispense pas d’étudier le domaine : les bornes d’un modèle peuvent représenter des situations limites.}\label{fig-bilan-A-cor-1}
\figureGeo[corrige]{bilan-A-aire}

\end{blocEdition}
\editionPage{bilan-A-cor-2}{Bilan A · Corrigé 2}{Deux médianes se rencontrent}{\hyperref[bilan-A-ex-2]{Énoncé p.\,\pageref*{bilan-A-ex-2}}}
\begin{blocEdition}{bilan-A-cor-2}
\editionEtape{1. Construire les milieux}
Les formules donnent $\res{I(5;3)}$ et $\res{J(1;3)}$.
\editionEtape{2. Déterminer les droites}
$(AI)$ a pour pente $3/5$ et passe par l’origine : $y=3x/5$. $(BJ)$ a pour pente $3/(1-8)=-3/7$ et passe par $B$ : $y=-3x/7+24/7$.
\editionEtape{3. Résoudre le système}
$3x/5=-3x/7+24/7$ donne $36x=120$, soit $x=10/3$, puis $y=2$. Donc $\res{G(10/3;2)}$.
\editionEtape{4. Contrôler la position}
$\vect{AG}=(10/3;2)=\tfrac23(5;3)=\tfrac23\vect{AI}$. Le point est sur la médiane et aux deux tiers de son trajet depuis $A$. Le résultat concorde avec le défi sur le point de concours des médianes.
\editionRetenir{Une équation de droite se vérifie sur deux points ; une intersection se vérifie dans les deux équations.}\label{fig-bilan-A-cor-2}
\figureGeo[corrige]{bilan-A-geometrie}

\end{blocEdition}
\editionPage{bilan-A-cor-3}{Bilan A · Corrigé 3}{Une enquête et son arbre}{\hyperref[bilan-A-ex-3]{Énoncé p.\,\pageref*{bilan-A-ex-3}}}
\begin{blocEdition}{bilan-A-cor-3}
\editionEtape{1. Reconstituer les cases}
Les effectifs des lignes sont $(48;32;80)$ et $(36;84;120)$. Les colonnes donnent $\res{84\text{ utilisateurs et }116\text{ non-utilisateurs}}$.
\editionEtape{2. Choisir les références}
La fréquence globale est $84/200=\res{0{,}42}$. Les conditionnements valent $P_A(S)=48/80=\res{0{,}6}$ et $P_B(S)=36/120=\res{0{,}3}$.
\editionEtape{3. Relier tableau et chemin}
$P(A)=80/200=0{,}4$ ; le chemin $A$ puis $S$ vaut $0{,}4\times0{,}6=\res{0{,}24}$, comme $48/200$.
\editionEtape{4. Inverser le conditionnement}
Parmi les 84 utilisateurs, 48 sont de $A$, donc $\res{P_S(A)=4/7\approx57{,}1\,\%}$. Le dénominateur 84 diffère du dénominateur 80 de $P_A(S)$.
\editionRetenir{Le tableau donne les effectifs ; l’arbre donne les probabilités de chemins. Les deux représentations doivent rester cohérentes.}\label{fig-bilan-A-cor-3}
\figureGeo[corrige]{bilan-A-arbre}

\end{blocEdition}
\editionPage{bilan-B-cor-auto}{Bilan B · Corrigé}{Automatismes expliqués}{\hyperref[bilan-B-auto]{Énoncé p.\,\pageref*{bilan-B-auto}}}
\begin{blocEdition}{bilan-B-cor-auto}
\begin{questions}
\item Diviser par 21 : $\res{2/3}$.
\item $2^{-2}=\res{1/4}$.
\item Différence de carrés : $\res{(x-3)(x+3)}$.
\item Un facteur doit s’annuler : $\res{S=\{-1;2\}}$.
\item Les membres sont positifs : $\res{x=9}$.
\item Le coefficient de $x$ est la pente : $\res{-3}$.
\item $2\times3-1\times4=\res2$.
\item $1{,}25\times0{,}80=1$ : taux global $\res{0\,\%}$.
\item La troisième valeur est $\res5$.
\item 2, 3, 4, 5 et 6 : $\res5$ valeurs.
\end{questions}
\end{blocEdition}
\editionPage{bilan-B-cor-1}{Bilan B · Corrigé 1}{Des durées regroupées}{\hyperref[bilan-B-ex-1]{Énoncé p.\,\pageref*{bilan-B-ex-1}}}
\begin{blocEdition}{bilan-B-cor-1}
\editionEtape{1. Cumuler}
$N=20$ ; les cumuls sont 4, 12 et 20, soit $\res{20\,\%,60\,\%,100\,\%}$.
\editionEtape{2. Estimer avec des hypothèses}
Les centres sont 5, 15 et 25 : la moyenne estimée vaut $(20+120+200)/20=\res{17\text{ min}}$. Elle serait exacte sous répartition uniforme dans les classes. La classe médiane est $[10;20[$ ; l’interpolation donne $10+(10-4)\times10/8=\res{17{,}5\text{ min}}$.
\editionEtape{3. Utiliser l’information exacte}
Avec les moyennes réelles, la somme vaut $4\times6+8\times14+8\times27=352$. La moyenne exacte est $\res{17{,}6\text{ min}}$.
\editionEtape{4. Ajouter une observation}
La nouvelle moyenne vaut $(352+60)/21=\res{412/21\approx19{,}6\text{ min}}$. La durée ajoutée est supérieure à l’ancienne moyenne, elle l’augmente. L’effectif est maintenant 21.
\editionRetenir{L’estimation par classes doit être remplacée par le calcul exact quand les moyennes de classes sont disponibles.}\label{fig-bilan-B-cor-1}
\figureGeo[corrige]{bilan-B-stats}

\end{blocEdition}
\editionPage{bilan-B-cor-2}{Bilan B · Corrigé 2}{Un rectangle et son cercle}{\hyperref[bilan-B-ex-2]{Énoncé p.\,\pageref*{bilan-B-ex-2}}}
\begin{blocEdition}{bilan-B-cor-2}
\editionEtape{1. Lire les directions}
$AB$ et $CD$ sont horizontaux, $AC$ et $BD$ verticaux. Les côtés consécutifs sont perpendiculaires : $ABDC$ est un rectangle de côtés 4 et 3, d’aire $\res{12}$.
\editionEtape{2. Calculer le centre}
Le milieu de $[BC]$ est $\res{O(3;5/2)}$. Pour chacun des quatre sommets, les écarts de coordonnées ont pour valeurs absolues 2 et $3/2$. Le carré de chaque distance à $O$ vaut $4+9/4=25/4$, donc la distance vaut $\res{5/2}$.
\editionEtape{3. Justifier le cercle}
Les quatre points sont donc sur le cercle de centre $O$ et rayon $\res{2{,}5}$. Les médiatrices de $[AB]$ et $[AC]$ sont $\res{x=3}$ et $\res{y=5/2}$ ; elles se coupent bien en $O$.
\editionRetenir{L’ordre des sommets du rectangle est indiqué : $A,B,D,C$. Un ordre mal choisi peut transformer un quadrilatère simple en un tracé croisé.}\label{fig-bilan-B-cor-2}
\figureGeo[corrige]{bilan-B-geometrie}

\end{blocEdition}
\editionPage{bilan-B-cor-3}{Bilan B · Corrigé 3}{Une urne sans remise}{\hyperref[bilan-B-ex-3]{Énoncé p.\,\pageref*{bilan-B-ex-3}}}
\begin{blocEdition}{bilan-B-cor-3}
\editionEtape{1. Décrire les changements de composition}
Au premier tirage, $P(R_1)=4/5$ et $P(B_1)=1/5$. Après rouge, il reste 3 rouges et 1 bleue : branches $3/4$ et $1/4$. Après bleu, il reste 4 rouges et aucune bleue : branches 1 et 0.
\editionEtape{2. Multiplier les chemins}
$P(R_1\cap R_2)=(4/5)(3/4)=\res{3/5}$. Les deux chemins de couleurs différentes valent chacun $\res{1/5}$ : $(4/5)(1/4)$ et $(1/5)\times1$.
\editionEtape{3. Utiliser une impossibilité}
Si la seconde est bleue, la première ne pouvait pas l’être car il n’y avait qu’une boule bleue. Elle est donc rouge avec certitude : $\res{P_{B_2}(R_1)=1}$. La branche de poids nul explicite l’impossibilité de deux bleues.
\editionRetenir{Sans remise, les probabilités du second tirage dépendent du premier. Une branche nulle peut apporter une information certaine après conditionnement.}\label{fig-bilan-B-cor-3}
\figureGeo[corrige]{bilan-B-arbre}

\end{blocEdition}
